% =====================================================================
%  sections/02_related_technique.tex  --  II. RELATED TECHNIQUE
%
%  Each entry keeps the same shape: a short definition, then a short
%  paragraph on why it was chosen for this project.
% =====================================================================
\section{RELATED TECHNOLOGIES}
\label{sec:related-technologies}

% ---------------------------------------------------------------------
\subsection{Core Technologies}
\label{sub:core-technologies}
This section outlines the core programming languages, frameworks, and technologies used to build \projname{}, along with the rationale behind each choice.

\subsubsection{Python}
Python is an interpreted, object-oriented, high-level programming
language known for its readability and its large ecosystem of libraries
for data analysis and machine learning.

\noindent \textbf{Why we chose it.} Every core AI and data component of
\projname{}, including the Retrieval-Augmented Generation (RAG) schema
pipeline, vector index management, and time-series forecasting engine, is
distributed as a Python library. Using Python across the backend unifies
\projname{}'s analytical and generative workflows into a single runtime,
eliminating inter-service serialization overhead and language boundaries.

\subsubsection{Large Language Models (LLMs)}
An LLM is an AI model trained on large volumes of text that can
recognise patterns, understand context, and produce coherent responses
across many language-based tasks.

\noindent \textbf{Why we chose it.} Translating user financial prompts into
executable PostgreSQL queries requires parsing natural language that is
frequently ambiguous, domain-specific, or variable in structure. An LLM
serves as \projname{}'s primary translation engine, handling diverse user
phrasings, schema relationships, and query formulations far beyond the
scope of static keyword parsers.

\subsubsection{Next.js}
Next.js is a React framework for building full-stack web applications,
with server-side rendering, incremental static regeneration, and
integrated API routes.

\noindent \textbf{Why we chose it.} The \projname{} interface requires reactive,
real-time updates across multiple components, including conversational chat,
financial dashboards, custom report builders, and forecasting modules.
Next.js provides a robust React component architecture and server-side
rendering that allows \projname{} to stream model responses and dynamic
charts seamlessly.

\subsubsection{Tailwind CSS}
Tailwind CSS is a utility-first CSS framework that styles an application
through utility classes instead of custom stylesheets, and keeps the
result responsive.

\noindent \textbf{Why we chose it.} With multiple team members collaborating
simultaneously on \projname{}'s user interface, Tailwind CSS enforced consistent
design standards across spacing, color schemes, and typography. Its
utility-first approach ensured responsive rendering for \projname{}'s data tables
and interactive visualization panels without CSS merge conflicts.

\subsubsection{FastAPI}
FastAPI is a high-performance Python web framework built on standard
type hints, using Pydantic for validation and generating OpenAPI
documentation automatically.

\noindent \textbf{Why we chose it.} FastAPI seamlessly connects \projname{}'s Next.js frontend
to the underlying Python AI pipeline. Its asynchronous execution handles
long-running LLM query processing and database fetches efficiently, while
automatic Pydantic data validation ensures structured API request payloads
between \projname{}'s components.

% ---------------------------------------------------------------------
\subsection{Libraries}
\label{sub:libraries}
This section outlines the core libraries and frameworks used to power \projname{}'s AI and data pipeline, along with the rationale behind each choice.

\subsubsection{LangChain}
LangChain is a framework for building applications powered by large
language models, providing abstractions that connect models to external
data sources, memory, and workflows.

\noindent \textbf{Why we chose it.} LangChain supplies the modular abstractions required for
\projname{}'s natural-language-to-SQL engine, specifically prompt templating,
vector database connection, and conversational chat memory. This allowed
the development effort to focus directly on schema retrieval optimization and
SQL generation logic.

\subsubsection{LangGraph}
LangGraph extends LangChain with a graph-based architecture, letting
developers express stateful, multi-step workflows as nodes and edges
with explicit control logic.

\noindent \textbf{Why we chose it.} Processing a user query in \projname{} involves a multi-step
decision flow rather than a simple prompt-response pair. LangGraph models
\projname{}'s execution graph, enabling deterministic routing such as short-circuiting
out-of-scope prompts, validating SQL syntax, and applying automated retry logic
if a database query fails.

\subsubsection{Chroma}
Chroma is an open-source, AI-native vector database that stores
embeddings and runs locally with minimal setup.

\noindent \textbf{Why we chose it.} Financial database schemas contain too many tables,
columns, and relationships to fit inside an LLM's prompt window. \projname{} uses
Chroma to index table schemas and retrieve only the most relevant schema fragments
for a given query, while keeping all data locally hosted for strict database privacy.

\subsubsection{Prophet}
Prophet is an open-source forecasting library from Meta (Facebook) for
time series data, modelling trend, seasonality, and holiday effects.

\noindent \textbf{Why we chose it.} Financial datasets like Kyungrinara exhibit strong
weekly, monthly, and seasonal purchasing cycles. Prophet powers \projname{}'s
forecasting engine by calculating time-series trends and generating upper and
lower uncertainty bounds, which \projname{} visualizes directly as confidence intervals
on its forecasting page.